% LTeX: language=en-GB
\chapter{Results: through the Workflow}
\label{chapter:results}

Do not include too many implementation details

\section{Mission and failure scenario analysis: calculating Vx and Vy}
* Mi servono degli input: prendo le informazioni raccolte nella fase preliminare di analisi missione (ref analisi missione)per stabilire probabilità di impatto and what not. 

* traiettoria
* punto di partenza, punto di arrivo
* drone scelto
* payload scelto
* medicinali scelti per simulazione

* serve energia di impatto?

---> risultati del WP2

\section{Studying the payload: calculating allowable impact acceleration}

\subsection{Material characterisation}
\subsubsection{TEST}
\subsubsection{FEM}
MENTION MODEL ITERATIONS TRIAL-AND-ERROR IN THE APPENDIX \ref{appendix:FEMiterations}!!!! FROM ABAQUS IMPLICIT TO CURRENT DECISIONS. Step by step, as in "version 1, abaqus implicit". Version 2, ls-dyna tetra elastic. Etc. MODAL PERTURBATION


* impact of container in LS-Dyna and what not. 
* diversi angoli di assetto per vedere il caso peggiore di accelerazioni locali
* cosa  possono sopportare le fialette e il contenitore????
* tipi di impatto, contenitore tipico che le contiene (menzionare la tesi del tipello) e che per il dimensionamento non viene considerato, approssimato a infinitamente rigido e che trasmette le accelerazioni di impatto così come le riceve ----> airbag deve essere dimensionato così
--------> tutte le analisi statistiche di prima, che hanno tenuto in considerazione anche il fatto che prima o poi si sorvolerà la città una volta risolti i cavilli burocratici che impongono una certa mitigazione del rischio ---> airbag. 

----> risultati del WP3

----> CONCLUSIONE: I have what I need to design the damn airbag. 


\subsection{Sled test}

\section{Design variables optimisation: the code!!!}\label{section:optimisation-code}

\begin{lstlisting}[language=Python, frame=single,
  literate={·}{{$\cdot$}}1 {→}{{$\rightarrow$}}2 {−}{{-}}1
           {ρ}{{$\rho$}}1 {φ}{{$\varphi$}}1 {ṁ}{{$\dot{m}$}}1
           {≤}{{$\leq$}}1 {≥}{{$\geq$}}1]
init: geometry(h=0), m0 = P0·V0/(R·T0)
loop while t < t_max:
  choose dt from previous load (hysteresis)
  F_n = max((P−Pamb)·A_prev, 0) → a_x, a_y          [explicit: old P, old A]
  u_new = u + a·dt ;  h_new = clamp(h + u_n,new·dt)  [semi-implicit Euler]
  A, V = geometry(h_new) ;  A_vent = A_or + φ(S − A)
  ρ = m/V → P, T adiabatic from previous state
  repeat N sub-steps (dt/N):                          [sub-cycling, not iteration]
      if P ≤ Pamb: stop
      dm = ṁ(P,T,A_vent)·dt/N → m, ρ, P (old T), T
  P = max(P, Pamb)
  if P > P_burst → BURST  | if h ≥ D0 → BOTTOMED | if y_pen ≤ 0 → REBOUND
  advance state
\end{lstlisting}

\begin{table}[h!]
\centering
\caption{Airbag design variables by design phase.}
\label{tab:airbag-design-variables-selected}
\begin{thesistable}{
  colspec = {X[l] X[l]},
  % ---- phase headers (blue = Phase 1, teal = Phase 2) and tinted contents
  row{1} = {c},
  cell{1}{1} = {bg=phaseA}, cell{1}{2} = {bg=phaseB},
  cell{2}{1} = {bg=tintA},  cell{2}{2} = {bg=tintB},
  row{2} = {abovesep=5pt, belowsep=6pt},
  vline{2} = {wd=\gap, fg=white},
}
Phase 1 -- design for peak acceleration & Phase 2 -- design for deployment time \\
\begin{minipage}[t]{\linewidth}
\begin{enumerate}[leftmargin=*, topsep=0pt, partopsep=0pt, parsep=0pt, itemsep=3pt,
                  label=\textcolor{phaseA}{\bfseries\arabic*.}]
  \item Shape
  \item Characteristic size
  \item Inflation pressure
  \item Fabric type, thickness
  \item Venting mechanism - vent burst pressure and/or fabric porosity
\end{enumerate}
\end{minipage}
&
\begin{minipage}[t]{\linewidth}
\begin{enumerate}[leftmargin=*, topsep=0pt, partopsep=0pt, parsep=0pt, itemsep=3pt,
                  label=\textcolor{phaseB}{\bfseries\arabic*.}]
  \item Housing/container size and geometry
  \item Air injection inlet pressure
  \item Air injection mass flow rate
  \item Seams locations
  \item Folding pattern
  \item Number of folds
\end{enumerate}
\end{minipage}
\\
\end{thesistable}
\end{table}

* codice ispirato da Zhou et al 2019 - spiegazione del modello e dell'aspetto dei vari blocchi
* Sapendo che le thresholds delle provette sono tot, allora si può convergere verso un pre-dimensionamento dell'airbag con tot dimensioni o caratteristiche. 
* oracle code for direzionamento verso un punto di partenza per il design

* I CHOSE THE CRUDE IMPLEMENTATION --- THIS IS BECAUSE THE ONLY COUPLING ANALYTICAL MODEL IS MISSING AND I CANNOT RETRIEVE IT. THEREFORE THIS IS MY IMPLEMENTATION XXX
Title: "Design, Fabrication and Testing of an Airdrop Platform Utilizing Airbags as Shock Absorbers"

Author: Timothy F. Patterson

Report: NATICK/TR-87/042 (DTIC accession ADA199593)

Date: August 1985 (work performed March 1984–May 1985)
\section{Summary of results}
Actual design variables p inflation, diameter, shape, etc

**RESULTS: BRING IT BACK TO THE RESEARCH QUESTIONS!!!!!**


